% Converted from the frozen Cycle026 Results draft; no new empirical evidence.
\paragraph{Counterfactual identity-memory evaluation.}
We hold the observation and semantic query fixed while switching the supplied
miner identity memory, which changes the intended helmet target.
Table~\ref{tab:primary-memory} reports the primary 50-group documented-protocol-disjoint
confirmation. CMA achieves 92\% CMSA on clean observations and 58\% under
\texttt{target15\_b}, compared with 4\% and 2\% for pinned SegLLM.
Under degradation, Memory Fidelity is 84\% versus 46\%, and Identity Error Rate
is 3\% versus 38\%. These results support stronger supplied-memory use on this
evaluated set. Figure~\ref{fig:main-memory} illustrates the controlled identity
switch and a retained degradation failure using previously frozen cases.

\paragraph{Robustness under compound degradation.}
CMA attains 67.53\% mIoU under the fixed compound stressor, compared with
36.44\% for SegLLM, a paired gap of 31.08 percentage points on identical groups.
The supported claim is higher absolute degraded performance and retained identity
fidelity: CMA still fails strict degraded CMSA on 21 of 50 groups.
The overlap-affected historical and development replications are reported
separately in Table~\ref{tab:overlap-replications} and are not pooled with the
primary set.

\paragraph{Comparison scope.}
SegLLM is a pinned released historical-memory comparator receiving the same
observations and supplied miner identities through its native interface.
The comparison is system-level: training exposure and domain adaptation differ,
and SegLLM's lower clean performance limits architecture-only attribution.
The counterfactual intervention measures identity-dependent behavior under common
inputs; it does not isolate the contribution of an individual CMA component.

\paragraph{Limitations.}
The primary set was selected before inference and is disjoint by source-image
SHA256 from reconstructed documented final-stage training and historical
evaluation registries. Exact run-bound w15 and ancestor/pretraining exposure
remain unknown, so this is not an exact training-unseen test.
Memories are externally supplied and targets are reconstructed pseudo labels;
23 groups retain historical review-level episode QC and 27 combine accepted-high
pairs. Structural validity does not replace independent annotation.
CMA's clean-to-degraded mIoU decline is 20.75 points, versus 6.26 for SegLLM;
we do not claim lower degradation sensitivity.
The evidence does not establish autonomous memory writing, persistence or repair,
broad corruption generalization, or safety guarantees. Earlier Layer-2
verifier/controller attempts did not meet their stated requirements and remain
negative evidence, not a successful agent contribution.
